\scanpage{27}
\uline{Cas particulier} Celui où $(xy)(zt)=0$ \inserted{[i.e. $(H^1)\deleted{\unsure{H|\otimes}}^{3}=0$ dans $H^4$ !]} quels que soient $x,y,z,t$, i.e. où l'image de $\alpha:\bigwedge^2 V\to k\xi+W$ est \deleted{\unsure{dans}} \deleted{\unsure{le|un}} ss-espace isotrope de $k(\xi)+W$. \deleted{Par exemple, si $W$ \unsure{a}, dans les alg. des\unsure{,} cela se réalise néc. ainsi : Si $W$ est \inserted{de dim} impaire, alors on peut \unsure{écrire}} $Q$ non dég. étant donnée, et supposant $W\ne0$, il existe (si $k$ alg. clos disons) une droite isotrope \inserted{$D$} dans $k\xi+W$ qui n'est pas dans $W$, et on prendra pour $\alpha$ \deleted{\unsure{une appl}} définie par une \deleted{forme}\inserted{\unsure{$\cdot W^L$|appl.}} bilinéaire \inserted{alt.} non dégénérée \deleted{sur} de $V\times V$ dans $D$ : alors $\varphi$ sera aussi non isotrope\ldots.

\begin{center}\rule{3em}{.4pt}\end{center}

Plus généralement, considérons la forme 4-linéaire alternée $\chi$ définie par
\[
xyzt=\chi(x,y,z,t)\xi^2\qquad (x,y,z,t\in V)
\]
Alors la forme \inserted{bilinéaire} induite par la \inserted{forme} $Q\oplus1$ (sur $\unsure{k|V}\xi+V$) sur $\bigwedge^{\unsure{2|?}}V$ par $\alpha$, qui définit une \deleted{forme} forme linéaire sur
\diagram{\begin{tikzcd}[ampersand replacement=\&,column sep=large]
\operatorname{Sym}^{2}(\bigwedge^2 V)\arrow[rr,bend right=18,"\text{\unsure{appl.|canon.}}"']\&\text{, se factorise en}\&\bigwedge^4 V
\end{tikzcd}}
moyennant $\chi$, et donc les données sont les suivantes

a) $V$ vect. (de dim finie) sur $k$

\deleted{b) $\varphi:\bigwedge^2 V\to k$}

b) $\chi:\bigwedge^4 V\to k$, définissant une forme bilinéaire symétrique $\chi'$ sur $\bigwedge^2 V$

c) un vectoriel $\widetilde W$, avec une forme \deleted{quadratique}\inserted{bil. sym.} $\widetilde Q$, et une application linéaire
\[
\bigwedge^2 V\xrightarrow{\alpha}\widetilde W
\]
compatible avec $\chi',Q$ [on dira que $\widetilde W,\widetilde Q$ est une \uline{extension} de la forme $\chi'$]

d) un élément $\xi$ de $\widetilde W$, tel que $Q(\xi,\xi)=1$.

\uline{Conditions de non \unsure{singularité|dégénérescence}}
\[
\left|\begin{aligned}
&Q\ \text{non dégénérée}\\
&\varphi(\unsure{x|s},y)=Q(\alpha(x\wedge y),\xi)\quad\text{non dégénérée}
\end{aligned}\right.
\]
% archive sheet number: 26
